% Source: https://artofproblemsolving.com/community/c3h2405045_tangency
% Stable identity: aops-topic-2405045-q01
Let $\ell$ be the line $O_1O_2$ and let $r_1,r_2$ be the radii, so $O_1T=r_1$ and $O_2T=r_2$.

Suppose, for contradiction, that $T$ is not on $\ell$. Let $T'$ be the reflection of $T$ in $\ell$; then $T'\ne T$. A reflection preserves distances and fixes every point of $\ell$, in particular $O_1$ and $O_2$. Hence
\[O_1T'=O_1T=r_1\qquad\text{and}\qquad O_2T'=O_2T=r_2,\]
so $T'$ lies on both circles. Then the circles have two common points $T$ and $T'$, contradicting the assumption that $T$ is the only one.

\begin{asy}
usepackage("amsmath");
size(5.5cm);
defaultpen(fontsize(10pt));
pair O1=(0,0), O2=(2.2,0); real r1=1.5, r2=1.2;
real xx=(r1^2-r2^2+2.2^2)/(2*2.2); real yy=sqrt(r1^2-xx^2);
pair T=(xx,yy), Tp=(xx,-yy);
draw(circle(O1,r1),blue); draw(circle(O2,r2),red);
draw((-1.9,0)--(3.8,0),dashed);
draw(O1--T--O2,gray); draw(O1--Tp--O2,gray);
draw(T--Tp,dotted);
dot(O1); dot(O2); dot(T); dot(Tp);
label("$O_1$",O1,SW); label("$O_2$",O2,SE); label("$T$",T,N); label("$T'$",Tp-(0,0.33));
label("$\ell$",(3.8,0),E);
\end{asy}

Therefore $T$ lies on $\ell$. $\blacksquare$

\emph{Remarks.} (1) The two cases are external tangency, with $T$ between the centres and $O_1O_2=r_1+r_2$, and internal tangency, with both centres on the same side of $T$ and $O_1O_2=\left|r_1-r_2\right|$.

\begin{asy}
usepackage("amsmath");
size(7cm);
defaultpen(fontsize(10pt));
pair A1=(0,0), A2=(1.8,0); real r1=1, r2=0.8;
draw(circle(A1,r1),blue); draw(circle(A2,r2),red);
draw((-1.3,0)--(2.9,0),dashed); draw((1,-1.1)--(1,1.1),gray);
dot(A1); dot(A2); dot((1,0));
label("$O_1$",A1,S); label("$O_2$",A2,S); label("$T$",(1.1,0.05),NE);
label("external",(0.8,-1.35));
pair B1=(4.75,0), B2=(5.52,0); real s1=1.45, s2=0.68;
draw(circle(B1,s1),blue); draw(circle(B2,s2),red);
draw((3.1,0)--(6.7,0),dashed); draw((6.2,-1.2)--(6.2,1.2),gray);
dot(B1); dot(B2); dot((6.2,0));
label("$O_1$",B1,SW); label("$O_2$",B2,S); label("$T$",(6.2,0),SE);
label("internal",(4.75,-1.8));
\end{asy}

(2) The common tangent line at $T$ is perpendicular to $O_1T$ and to $O_2T$ (a tangent is perpendicular to the radius at the point of contact), which is another way to see that $O_1$, $T$, $O_2$ are collinear once one knows that tangent circles share a tangent line at $T$.
